\section{Proofs}
\label{app:proofs}

We restate and prove the two propositions of Section~\ref{sec:identif}. Throughout,
$P$ denotes the population law of the (weakly stationary) process, and the
one-step reconstruction risk of a predictor $g=(g_1,\dots,g_d)$ is
$R(g)=\sum_{i=1}^{d}\mathbb{E}_P\big[(x_{t,i}-g_i)^2\big]$.

\subsection*{A.1\quad Identifiability of SPADE's target (Proposition~\ref{prop:identif})}

\begin{proof}[Proof sketch]
The argument has three steps.

\emph{Step 1 (SPADE's target is a bona fide restricted ANM).} By construction
(Eq.~\ref{eq:predict}), SPADE predicts each target through additive univariate
functions of its candidate parents and nothing else, so any element of its
hypothesis class has the form $x_{t,j}=\sum_{h,i} f^{(h)}_{ij}(x_{t-h,i})+e_{t,j}$
with mutually independent residuals---identical in form to the data-generating
Eq.~\eqref{eq:svar}. Under Assumption~\ref{as:markov} (causal minimality) each
mechanism depends non-trivially on its argument, so the model is a
\emph{restricted} additive-noise model in the sense of
\citet[\S 4]{peters2014causal}.

\emph{Step 2 (lagged and contemporaneous identification).}
For lagged edges, condition on the past: given
$\mathbf{x}_{<t}$, the map $\mathbf{x}_{t-h}\mapsto x_{t,j}$ is a regression whose
non-vanishing partial dependence on coordinate $i$ certifies
$A^{(h)}_{ij}=1$; time ordering fixes the direction, so
$\{A^{(h)}\}_{h\ge 1}$ is identified from the joint law without an acyclicity
requirement. For the contemporaneous block, restrict to the instantaneous SEM
over $\mathbf{x}_t$: by \citet[Thm.~31]{peters2014causal}, a restricted ANM
satisfying causal minimality has its DAG identifiable from the observational
distribution, provided the degenerate configuration of
\citet[Cond.~18]{peters2014causal} does not hold. Assumption~\ref{as:restrict}
excludes exactly that configuration (of which linear--Gaussian is the canonical
instance, cf.\ \citealt{hoyer2008nonlinear}), giving unique identification of
$A^{(0)}$; without it, $A^{(0)}$ is identified only up to Markov equivalence, the
general statement in part~(b).

\emph{Step 3 (the constrained population minimizer recovers $A$).} B-splines of
order $p$ are dense in the H\"older/Sobolev classes as the grid refines, and KAN
regressors built from them attain the minimax-optimal nonparametric rate
\citep{Liu2024,liu2025kanconvergence}; hence for a sufficiently fine grid the
true mechanisms lie in (the closure of) SPADE's class and the population risk
$R$ is minimized at the true regression functions. Because, by Steps~1--2, the
true functions correspond to a \emph{unique} adjacency once the feasible set is
restricted to DAGs on the contemporaneous block ($h(W^{(0)})=0$,
Theorem~\ref{thm:notears}), the constrained minimizer's summary graph equals $A$.
Finite-sample recovery is not claimed; the statement is a well-specified,
large-sample identifiability guarantee, consistent with our empirical finding
that recovery is strongest where the target is genuinely identifiable (non-linear
instantaneous structure) and weaker on near-integrated lagged series.
\end{proof}

\subsection*{A.2\quad Necessity of the information bottleneck (Proposition~\ref{prop:bottleneck})}

\begin{proof}
Let the predictor for target $i$ be $g_i=\psi_i(\mathbf{x}_{-i})$ drawn from a
class $\mathcal{F}$ of functions of all other coordinates
$\mathbf{x}_{-i}$, and suppose $\mathcal{F}$ is rich enough to contain the
conditional mean $m_i(\mathbf{x}_{-i})=\mathbb{E}_P[x_{t,i}\mid\mathbf{x}_{-i}]$.
The per-target risk $\mathbb{E}_P[(x_{t,i}-\psi_i)^2]$ is minimized, over all
measurable $\psi_i$, uniquely (in $L^2(P)$) at $\psi_i=m_i$, and its value depends
only on $\psi_i$. Now suppose an adjacency $\hat A$ is read \emph{post hoc} from
the fitted predictor by any rule that does not itself enter the objective. Since
the objective $R(g)=\sum_i\mathbb{E}_P[(x_{t,i}-\psi_i)^2]$ contains no term in
$\hat A$, its minimizer $g^\star=(m_1,\dots,m_d)$ is unchanged for every choice of
$\hat A$. Consequently all adjacencies are risk-equivalent at the optimum and
$\hat A$ is not identified by risk minimization.

Conversely, impose the bottleneck: require
$g_i=\sum_{j}\mathbf{1}[\hat A_{ji}=1]\,\phi_{ji}(x_j)$, i.e.\ $\hat A$ indexes
\emph{which} additive edge terms are permitted. Then the attainable function set
is $\mathcal{G}(\hat A)=\{\sum_{j\in\pa(i)}\phi_{ji}\}$, and $R$ becomes a genuine
function of $\hat A$ through the constraint $\mathcal{G}(\hat A)$. Under
Assumptions~\ref{as:markov}--\ref{as:restrict} the conditional mean $m_i$ is
additive in the true parents (Eq.~\ref{eq:svar}), so $m_i\in\mathcal{G}(\hat A)$
iff $\hat A$ contains the true parent set of $i$; among such $\hat A$, causal
minimality makes the true (minimal) parent set the unique risk minimizer once a
sparsity preference breaks ties toward fewer edges. Hence the risk-minimizing
structure is unique, which is what the bottleneck buys.
\end{proof}

\begin{remark}
Proposition~\ref{prop:bottleneck} is not specific to KANs: it explains why
\emph{any} sufficiently expressive shared-backbone predictor fails as a structure
learner, and why component-wise neural-Granger designs
\citep{Tank2021,GCKAN2024} and per-node score-based learners
\citep{zheng2020sparse,lachapelle2020grandag} adopt the same restriction. SPADE's
contribution is to combine this bottleneck with \emph{learnable spline edges},
which is what lets it identify the non-linear mechanisms that fixed-architecture
learners approximate poorly (Section~\ref{sec:instdag}).
\end{remark}

\section{Reproducibility}
\label{app:repro}

All experiments use the released code and fixed seeds. Unless stated otherwise,
SPADE uses $E{=}150$ epochs (400 for the instantaneous-DAG study), learning rate
$5\times10^{-3}$, spline grid $G{=}8$ of cubic ($p{=}3$) B-splines on
$[-3,3]$, group-lasso weight $\lambda_g{=}10^{-2}$ (headline instantaneous-DAG
study: $\lambda_g{=}10^{-2}$, selected via a sensitivity sweep on
held-out validation seeds (100--104/100--102) disjoint from the seeds reported in
Table~\ref{tab:instdag} (0--4/0--2), after finding the earlier $\lambda_g{=}0.02$ was
over-regularized for this task---see Section~\ref{sec:instdag}. An
intermediate value, $\lambda_g{=}5\times10^{-3}$, was reported in an earlier
revision but was selected by a tuning script that, we later found on
independent review, had erroneously swept the same seeds used to report
Table~\ref{tab:instdag} rather than the disjoint validation seeds its own
comment claimed; \texttt{scripts/tune\_instdag\_lambda.py} now uses the
disjoint seeds it always claimed to, and the value above is what that
corrected script selects. Because the initial disjoint sweep used only five
held-out seeds, we then expanded it to twenty (seeds 100--119) to check
whether $\lambda_g{=}0.01$ was a confident choice or an artifact of a small
sample: AUROC across $\lambda_g\in\{0.005,0.008,0.01,0.012,0.015,0.02\}$ was
$0.955/0.969/0.969/0.959/0.948/0.940$ respectively, with a paired Wilcoxon
test of $\lambda_g=0.01$ against $\lambda_g=5\times10^{-3}$ over the same
twenty seeds giving $p=0.39$---not a statistically confident difference.
$\lambda_g=0.01$ sits inside this broad, statistically indistinguishable
plateau (its peak is at $0.008$, itself not distinguishable from $0.01$),
while $\lambda_g=0.02$ is clearly outside it; we therefore keep
$\lambda_g=0.01$, the value already applied and reported throughout this
paper, rather than make a further change for a difference too small to
support one), gradient clipping at $5.0$,
and---when contemporaneous edges are modelled---an Augmented-Lagrangian schedule
$(\rho_0,\alpha_0,\beta,K)=(1,0,2,50)$ with $\rho$ capped at $10^{10}$. Series are
$z$-scored per training window (Remark~\ref{rem:varsort}). Baselines use their
authors' reference implementations: DAGMA \citep{Bello2022} via the
\texttt{dagma} package; GraN-DAG \citep{lachapelle2020grandag}, GOLEM
\citep{ng2020golem}, and non-linear NOTEARS \citep{zheng2020sparse} via
\texttt{gcastle}; PCMCI \citep{Runge2019} via \texttt{tigramite}; VAR-LiNGAM
\citep{Hyvarinen2010} via \texttt{lingam}; SCORE \citep{rolland2022score} and
NoGAM \citep{montagna2023nogam} via \texttt{dodiscover} (the py-why project's
\texttt{dodiscover.toporder} module, which the NoGAM authors' own repository
identifies as its reference release; not distributed on PyPI under that name,
so installed from a source clone of the project rather than \texttt{pip
install dodiscover}). Any baseline that failed to run in our
environment within the compute budget is reported as such rather than omitted.
NoGAM's kernel-ridge cross-validated pruning step did not complete within
budget for the full five-seed protocol at $d{=}20$; we report it there over a
single seed, flagged in Table~\ref{tab:instdag}, rather than omit the width or
extrapolate.

Three further settings were corrected during this revision, each via a
held-out procedure that never touched the seeds used for a reported number
before the setting was fixed. First, SPADE's forecasting head
(Section~\ref{sec:forecast}) originally trained for only 30 epochs against
the deep baselines' 80; a sweep on an inner train/validation split of the
training fold (never the test folds reported in Table~\ref{tab:forecast})
found validation error still improving through roughly 300 epochs, and we
adopted that as the epoch count for all further forecasting runs. Second,
the lagged non-linear causal-discovery benchmark
(Section~\ref{sec:synthetic}) uses a data generator
(\texttt{generate\_nonlinear\_scm}) whose lagged summary graph can contain
cycles and whose mechanisms include an unbounded term; a minority of
$(d,\text{seed})$ draws at $d\in\{10,20,50\}$ diverged to non-finite or
extreme-scale data, which classical baselines failed loudly on (and were
silently dropped by the benchmark's error handling) while gradient-based
methods, including SPADE, trained on silently without producing a finite,
informative fit. We added a finite-and-bounded check with deterministic
reseeding, scoped to \texttt{scripts/honest\_causal\_benchmark.py} only, and
reran the affected configurations once. Third, in the same benchmark's
\emph{linear} configuration at $d{=}50$, the group-lasso weight
$\lambda_g{=}0.01$ (which works well at $d\le20$) proved severely
over-regularized: the penalty sums a per-edge norm over all
$d(d{-}1)\times\text{max\_lag}$ candidate edges, so the same $\lambda_g$
grows relatively stronger as $d$ increases. A sweep on held-out validation
seeds found $\lambda_g\le0.003$ recovers AUROC$=1.0$ at this width, and we
adopted $\lambda_g{=}0.002$ (width-dependent choice documented in
\texttt{fit\_cdkan}'s docstring), rerunning only the $d{=}50$ rows on the
official seeds.
